
When LLM-generated code fails a test, one approach to repair involves providing the model with error feedback and requesting a fix. Prior work on self-debugging and iterative code repair has used various forms of error feedback, ranging from simple pass/fail signals to detailed execution traces with variable states. However, the effect of feedback granularity on repair success has not been systematically studied.

This work addresses the following research question: Does the granularity of error feedback affect LLM code repair success, and if so, how?

\subsubsection{1.1 Background}

The self-debugging paradigm, established by Chen et al. (2023), demonstrated that LLMs can improve their own code using execution feedback, achieving reported improvements on benchmarks such as MBPP. Subsequent work has extended this approach with execution traces \citep{bouzenia2024repairagent}, reinforcement learning from execution feedback \citep{gehring2024rlef}, and multi-agent frameworks \citep{Huang2026}. These approaches have used different levels of error detail, but none have systematically compared granularity levels as the primary variable.

\subsubsection{1.2 Research Gap}

Prior work varies feedback granularity incidentally rather than systematically. Self-Debug uses error messages (approximately G2-level feedback), TraceFixer uses full stack traces (G4-level), and DynaFix adds variable states. No prior controlled comparison exists that isolates granularity as the sole independent variable.

\subsubsection{1.3 Hypotheses}

This study tested the following hypotheses:

\begin{itemize}
\item \textbf{H-E1 (Foundation)}: Runtime errors with localizable stack traces constitute at least 30\% of LLM code generation failures on MBPP.
\item \textbf{H-M1 (Mechanism)}: Error feedback granularity has a statistically significant effect on repair success rate (ANOVA p < 0.05).
\item \textbf{H-M2 (Direction)}: G3 (error with line number) achieves at least 10 percentage points higher repair success than G0 (pass/fail only).
\item \textbf{H-M3 (Non-monotonicity)}: G4 (full stack trace) does not significantly outperform G3 (difference within 2\%).
\end{itemize}

The directional hypotheses (H-M2, H-M3) were based on an ''attention window hypothesis'' predicting that intermediate granularity would focus model attention optimally.

\subsubsection{1.4 Summary of Findings}

H-E1 and H-M1 were supported by the data. H-M2 and H-M3 were not supported; the observed pattern was opposite to the predicted direction for H-M2, and G4 significantly outperformed G3 contrary to H-M3.
